\documentclass[conference]{IEEEtran}
\usepackage{amsmath,amssymb,graphicx}

% Replace every bracketed placeholder before submitting.
% Compile with: pdflatex report_template.tex
% The suggested length is two pages; length itself is not graded.

\title{Slalom Challenge: [Your Method Name]}
\author{
  \IEEEauthorblockN{[Your Name]}
  \IEEEauthorblockA{[Student ID] \\
  Learning-Based Control for Mobility Systems, 2026}
}

\begin{document}
\maketitle

\section{Problem Statement}
% State the mode (easy or full), state and action variables, constraints,
% official minimum-time goal, and any surrogate training objective.
[Describe the slalom task and the problem your method actually solves.
Explain how your training objective relates to the official minimum-time
evaluation.]

\section{Proposed Approach}
% Explain the course concepts you implemented, not just an algorithm name.
% Give model/data sources, learning or optimization steps, offline outputs,
% and the online decision rule. Replace the placeholder below with an actual
% representative figure of your method or controller.
[Describe your controller and learning or optimization procedure.]

\begin{figure}[t]
\centering
\fbox{\begin{minipage}[c][0.75in][c]{0.88\columnwidth}
\centering
Replace with a representative method or controller figure.
\end{minipage}}
% Example replacement:
% \includegraphics[width=0.94\columnwidth]{method_figure.pdf}
\caption{[Explain what the figure shows and identify its main signals.]}
\label{fig:method}
\end{figure}

\section{Own Validation}
% Use results you measured. Include run settings, completion/failure,
% finish time if successful, a meaningful baseline, and limitations.
[Present your PyChrono results and comparison. Explain what the evidence
does and does not establish.]

\section{Conclusion}
[State the main result and the most important limitation.]

\noindent\textbf{AI use:} [List tool and model names and their main uses,
or state \emph{No generative AI used}.]

% Cite external code/data if used. The bibliography may be omitted if none.
% \begin{thebibliography}{1}
% \bibitem{source} Author, Title, venue or URL, year.
% \end{thebibliography}

\end{document}
